% Full ordinary mixed-term proof; no left-linearity or termination assumed.
\begin{annexitem}{annex:p12-klop1987-ds-invariants}{thm:p12-klop1987-ds-invariants}{Every new colour path comes from an old variable instance}
\begin{proof}
A rule of colour \(a\) replaces a monochromatic left pattern
by a monochromatic right pattern, or a captured variable instance.
Every variable on the right already occurs on the left. Each new
root-to-leaf path therefore consists of the old outside prefix,
a block of colour \(a\), and a path in an old captured instance.
The old left side already provided such a colour-\(a\) block
above that instance. Copying makes more copies of these paths,
but cannot add a colour change on any of them. A collapsing right
side removes the block and can merge adjacent equal colours.
Thus the maximum number of colour blocks, with minimum rank1,
cannot increase.

A pure one-signature term remains pure: matching in it assigns
variables only to its existing subterms, and the rule cannot
introduce the other signature or new free variables. It therefore
remains preserved, including when it becomes a variable.
For a mixed preserved term, prove closure by induction on its rank.
Its universal root condition remains true at every reduct.
An inner step changes one principal alien; the lower-rank induction
preserves that alien, and its root cannot enter the parent's colour:
a mixed preserved alien retains its initial colour, while a pure
opposite-colour alien can have only that colour or a variable.
An outer step changes only the cap and copies or erases captured
aliens. Since the whole root cannot change colour by its universal
guard, the resulting principal aliens are retained old aliens or
their preserved descendants. Thus the recursive preservation
conditions persist. This proves closure and the foreign-root
assertion. No step-count argument or global termination was used.
\end{proof}
\end{annexitem}

\begin{annexitem}{annex:p12-klop1987-ds-coherence}{thm:p12-klop1987-ds-coherence}{Finite join classes and the substitution that preserves every old equality}
\begin{proof}
For terms with CR at each of them, joinability is transitive on
their finite list: if \(A_i\downarrow A_j\) and
\(A_j\downarrow A_k\), their two common reducts are descendants
of \(A_j\); CR there joins them, giving a common reduct of
\(A_i,A_k\). Reflexivity and symmetry are immediate. The list
thus partitions into joinability classes. Every descendant retains
local CR, and joinability of two originals persists for any chosen
descendants of them: join those descendants with the old common
reduct, then join the resulting reducts there. Iterating this
finite joining argument in each class gives one common term below
all the finitely many chosen descendants in that class. Assign it
to every index in the class. Equal originals or equal descendants
are necessarily in the same class. The resulting tuple is therefore
coherent with every chosen tuple, as well as a common descendant
of all of them. Neither normal forms nor SN are required.

For transport, replace every distinct current frontier value by a
different fresh variable in the cap. A pattern of the cap's colour
cannot inspect the inside of a foreign-root or neutral frontier
value; it can only capture it at a variable. Literal equality of
captured cap subterms is reflected by this abstraction: their cap
operators and all repeated frontier values must agree. Thus the
actual outer step is an ordinary component step on the abstract cap.
Coherence makes the update a well-defined substitution on those
fresh variables. Applying it to the whole abstract step preserves
that step under substitution. Interpreting the updated variables
as their actual values lifts it back to the concrete cap.
Every right-side occurrence of a captured variable uses the same
updated value; multiple occurrences copy it, absent occurrences
erase it, and a variable right side selects it. This gives exactly
the stated index transport, without assuming left-linearity.

For holes at the deepest deletion frontier the same argument
applies to any redex properly above them. A path from that redex
to such a special-layer root crosses a signature boundary, so a
single-colour operator pattern cannot enter it. Equal copies of
one matched variable also have identical selected-layer positions:
their entries from the rule's monochromatic pattern have the same
layer depth. Replacing equal frontier values by equal ones therefore
preserves equality of the entire matched variable instances, not
only a selected pair of alien leaves.
\end{proof}
\end{annexitem}

\begin{annexitem}{annex:p12-klop1987-ds-preserved}{thm:p12-klop1987-ds-preserved}{The repaired-inner/outer block diamond and rank induction}
\begin{proof}
Consider \(H[\mathbf A]\) with the stated preserved frontier values
outside its cap colour \(a\), each with local CR. Their descendants
remain outside that colour by the invariant lemma. Every step is
therefore either inside a frontier value, or an \(R_a\)-step in
the cap. An outer step only changes the cap and copies, erases or
selects the frontier occurrences. If it selects one whole value,
regard the cap as the single hole thereafter; all further actual
steps are inside it. This keeps the family closed even under collapse.

Fix a tuple of frontier values. Abstract its distinct values by
distinct fresh variables as in the coherence lemma. Every outer
path is a component path, and every component joining path lifts
back under the interpretation substitution. CR of \(R_a\)
therefore proves confluence of outer-only paths at that fixed tuple.

Now take two paths consisting of an inner block followed by an
outer block. The initial inner blocks yield tuples \(\mathbf B\)
and \(\mathbf C\) from \(\mathbf A\). Synchronize them to
\(\mathbf D\) by the finite coherence lemma. Transport both
outer paths to the cap containing \(\mathbf D\).
Their original endpoints reduce to the transported endpoints by
reducing every retained or copied frontier value to its assigned
\(D_i\). Erased values require no step. The two transported
outer paths start at the same cap and the same tuple, so outer-only
confluence joins them. Both original endpoints consequently join
by another inner block and outer block.

Thus the relation consisting of an inner block followed by an
outer block has the diamond property. It contains every one-step
reduction of this family and is contained in ordinary finite
reduction. Its transitive closure is therefore the full finite
reduction of the family. A finite grid of its diamonds joins any
two finite block paths, proving CR of the original term. All tuple
CR premises persist after copying and reduction, so the same block
argument applies at every later corner of that finite grid.

Finally induct on layer rank for a preserved term. Pure terms have
CR by the corresponding component assumption; a variable has no
step. In a mixed preserved term, every principal alien is preserved
and has lower rank. The induction gives its CR. Ordinary variable
leaves are also preserved and CR. The cap argument now proves CR
of the whole term. If the root term itself is nonpreserved but all
its principal aliens are preserved, exactly the same cap argument
applies: it did not require preservation of the cap root, and
explicitly included its collapse. This proves the full stated lemma.
\end{proof}
\end{annexitem}

\begin{annexitem}{annex:p12-klop1987-ds-depth}{thm:p12-klop1987-ds-depth}{Finite semantic depth, its monotonicity and the actual deletion witnesses}
\begin{proof}
The principal-alien recursion is on proper finite subterms, so
\(d\) is a finite natural number, bounded by layer rank.
If a term is not preserved, descend through a nonpreserved principal
alien while one exists. This finite descent ends at a nonpreserved
layer with all aliens preserved. Its failure of the root guard gives
an actual finite reduction whose root leaves its initial signature.
Stop at the first such change. Before it, the monochromatic cap
has only preserved foreign/neutral frontier values; its first root
collapse selects a descendant of one of them. That term is preserved.
The cap lemma also proves CR of this entire deletion occurrence.

Two minimal nonpreserved special occurrences cannot contain each
other, since every special layer below either lies in its preserved
principal aliens. Their witness reductions can therefore be performed
sequentially in disjoint contexts. This proves that parallel deletion
is an ordinary finite reduction and is available whenever \(d>0\).

For the depth bounds, first observe that a one-colour cap applied
to any finite tuple of terms has depth at most one plus their
maximum depth. A minimal nonpreserved layer is either the cap's
own layer, or occurs within one tuple member; wrapping in one
colour adds at most one layer, while merging equal colours removes
one. This also applies when the cap is a single hole or has no
holes, taking the empty maximum as0.

Prove ordinary depth nonincrease by induction on \(d(t)\).
At0, preserved closure applies. At1, all principal aliens are
preserved; their descendants retain this status and cannot enter
the cap colour. Every later term has at most that one potentially
nonpreserved cap layer, or is a selected preserved value, so its
depth is at most1. At larger depth write
\(t=C[A_1,\ldots,A_k]\), with
\(d(t)=1+\max_i d(A_i)\). An inner step changes one \(A_i\);
the induction bounds its new depth, and the cap bound bounds the
whole result, including an alien that now merges into the cap's
colour. An outer step has only retained, copied or erased old
frontier values in its new one-colour cap, so the same bound applies.
A collapsed result has no extra cap layer. These cover all positions.

In a parallel deletion, all deepest minimal nonpreserved layers
become preserved. Their root changes remove their coloured layer
or merge it with the parent layer. No preserved region can become
nonpreserved under these actual reductions. Every surviving or
newly minimal nonpreserved layer is therefore either an unchanged
shallower one, or an ancestor of a deleted deepest layer, at strictly
smaller layer depth. Consequently \(d\) drops strictly. At depth1
the whole result is preserved and has depth0. This is termination
of the auxiliary deletion relation only, not of ordinary rewriting.
\end{proof}
\end{annexitem}

\begin{annexitem}{annex:p12-klop1987-ds-diagrams}{thm:p12-klop1987-ds-diagrams}{Disjoint, inside, and above the deepest frontier}
\begin{proof}
Let \(M_1,\ldots,M_k\) be the disjoint deepest deletion occurrences
of \(t\), and \(P_i\) their chosen root-changing reducts in
\(v\). Each \(M_i\) has CR and each \(P_i\) is preserved.
Compare the ordinary redex position with these occurrences.

If it is disjoint from all of them, their occurrences and witnesses
commute literally with the step. No new equally deep bad occurrence
can arise in the changed disjoint region: its old bad depth was
smaller and ordinary depth cannot increase. In the equal-depth
case perform the same parallel deletion in \(u\), then the
unchanged ordinary step in \(v\). The results coincide and
have depth below \(d(t)\). A depth drop cannot remove these
unaffected deepest occurrences, so the strict-drop case does not
occur at this position.

If the ordinary step is inside \(M_r\), let its new value be
\(N_r\). All its aliens started preserved, so its reduct still
belongs to that stable-frontier cap family. It is either preserved
or still a minimal nonpreserved layer. In the equal-depth case
delete it in \(u\) if it remains a deepest bad occurrence;
otherwise leave it, and use the unchanged witnesses on the other
deepest occurrences. Call its resulting value \(Q_r\).
Both \(P_r\) and \(Q_r\) descend from the CR term \(M_r\),
so join them at \(S_r\). Insert that join and all other chosen
\(P_i\) into their context. This gives the required common
\(z\), a descendant of \(v\), hence of smaller depth. If
the ordinary step already drops the global depth, take the
ordinary branch unchanged and join \(P_r,N_r\) in exactly
the same manner. There are then no other unaffected deepest bad
occurrences; otherwise global depth could not have dropped.

In the remaining case the ordinary redex lies above some selected
\(M_i\). Its monochromatic pattern cannot enter those special
layer roots, because the path crosses a signature boundary.
They occur inside its matched variable instances. On the ordinary
branch they are retained, copied, erased or moved by collapse.
For equal global depth, perform parallel deletion at the deepest
occurrences of \(u\). A copied \(M_i\) that moved to a
shallower depth need not be selected there: retain it unchanged
for the subsequent join. A selected copy has some chosen reduct
\(Q_{i,j}\), and an unselected copy has \(Q_{i,j}=M_i\).
Include all these finite choices and the old \(P_i\) in one
coherent synchronization below their original CR terms. It gives
\(D_i\) with \(P_i\to^*D_i\), every
\(Q_{i,j}\to^*D_i\), and \(M_i=M_l\Rightarrow D_i=D_l\).

Reduce the old deletion branch's values to \(D_i\). The exact
coherence transport lemma now restores the original rule match,
including every repeated variable, so perform that same outer
rule step. The new branch reduces each retained/copy value to
its \(D_i\); both endpoints are literally the same term \(z\).
Erased indices need no corresponding right-side step. Values
outside the redex are synchronized in their unchanged holes as
well. Since \(z\) descends from \(v\), its depth is below
\(d(t)\). If the ordinary step already lowered the depth, omit
the new parallel deletion and put \(Q_{i,j}=M_i\) for all
surviving copies. The same synchronization and transport gives
the join directly with \(u\).

These positions exhaust a redex relative to a finite disjoint
frontier. In particular no assumption that unequal copies reduce
identically was made; the finite CR synchronization establishes
the equality needed for each nonlinear match.
\end{proof}
\end{annexitem}

\begin{annexitem}{annex:p12-klop1987-ds-confluence}{thm:p12-klop1987-ds-confluence}{The depth/path double induction on all mixed terms}
\begin{proof}
Assume both components CR. Induct on \(d(t)\) to prove CR at
every term. At0 it is a preserved term and the cap/rank theorem
applies. Suppose CR is established for every depth below \(n>0\).
Prove first, by induction on path length \(m\), simultaneously
for all sources of depth \(n\), that a finite path
\(t\to^m u\) joins any chosen \(t\Rightarrow_{
m d}v\)
at a descendant of \(v\) having depth below \(n\).
For \(m=0\), use \(v\) itself, since the deletion is a finite
ordinary reduction from \(u=t\).

For \(m>0\), write \(t\to a\to^{m-1}u\).
If \(d(a)<n\), the strict-drop diagram gives a common reduct
\(b\) of \(a,v\), of depth below \(n\).
The induction on depth gives CR at \(a\), joining its paths
to \(u\) and \(b\). Their common reduct also descends from
\(v\), and its depth is below \(n\).
If \(d(a)=n\), the equal-depth diagram gives a deletion
\(a\Rightarrow_{
m d}w\) and a common reduct \(b\) of
\(v,w\), below depth \(n\). The shorter-path induction at
\(a\) joins \(u,w\) at \(c\), also below \(n\).
Since \(d(w)<n\), CR at \(w\) joins its reducts \(b,c\)
at \(e\). Then \(v\to^*b\to^*e\) and
\(u\to^*c\to^*e\). This proves the path claim without
using confluence at the still-unproved depth \(n\).

Finally take any two finite reductions from \(t\) and fix one
parallel deletion \(t\Rightarrow_{
m d}v\).
The path claim joins each endpoint with a reduct of \(v\).
Since \(d(v)<n\), confluence at \(v\) joins those two reducts.
Concatenation joins the original endpoints. This completes the
depth induction and proves confluence of the whole sum.

Conversely each pure component term set is successor-closed in
the sum: the other signature has no redex there, and a component
rule with right-side variables already on its left cannot introduce
a foreign operator into a matched pure term. Joining reductions
from pure endpoints therefore stay pure. CR of the sum supplies
CR of each component. Renaming signatures transports this argument
through the specified isomorphisms. No rule-linearity or SN premise
was inserted in either direction.
\end{proof}
\end{annexitem}

\begin{annexitem}{annex:p12-klop1987-ds-first-counter}{thm:p12-klop1987-ds-first-counter}{The two terminating factors and the exact sum cycle}
\begin{proof}
In the first factor a subterm can become0 or1 at its root only
if it was that constant initially: an f-root remains an f-root,
variables are fixed, and both constants are irreducible.
Thus an inactive f-node cannot acquire the first two arguments
0 and1. After a root contraction, its three copies of x cannot
later provide0 and1 simultaneously: all their possible root
categories agree, and only the constant categories can reach those
constants. The new root is permanently inactive.
Give every initial active redex an origin and its depth in the
finite containment forest of active redexes. All later active
redexes are copied descendants of these originals. Contracting
one origin removes one residual at its depth and can copy only
strictly deeper origins; shallower counts are unchanged. The
finite lexicographic tuple of residual counts, shallow first,
therefore decreases at every step. Lexicographic order on a
fixed finite product of naturals is well-founded, proving SN.

Each or-step selects a proper subterm, strictly decreasing node
count, so the second factor is SN. It is not CR: or(0,1) reduces
to the distinct irreducible constants0 and1. In the sum select
0 in the first q-copy,1 in the second, and apply the f-rule at
the resulting root. This is exactly the displayed three-step
cycle, and repetition is infinite. The foreign collapsing
operator makes the copied argument admit two different root
constants, which the first factor's pure-term argument excluded.
\end{proof}
\end{annexitem}

\begin{annexitem}{annex:p12-klop1987-ds-complete-counter}{thm:p12-klop1987-ds-complete-counter}{Completeness of both exact factors and their periodic mixed term}
\begin{proof}
For the first factor, the constant descendant sets of1,2,3 miss
at least one of4,5,6: they are contained respectively in
\(\{1,4,5,7\},\{2,4,6,7\},\{3,5,6,7\}\).
The sets for4,5,6,7 are smaller and likewise cannot supply all
three. An F-root can only remain F or become7, and a variable
cannot change. Thus, after a duplicating step
\(F(4,5,6,x)\to F(x,x,x,x)\), the resulting root can never
again match4,5,6, regardless of how its copies reduce separately.

Trace every F-node to its initial origin, tagging a root as
spent after its duplicating step and copying the tags with subterms.
Every duplicating step is at an unspent root by the preceding
argument. A copied proper descendant origin has greater initial
containment depth than that root. Use the finite tuple counting
unspent F-residuals by these depths, followed by the total number
of F-nodes and the sum of constant weights. Assign weights2 to
1,2,3;1 to4,5,6;0 to7. A duplicating step strictly lowers its
first affected unspent-origin count and can increase only later
coordinates. The catch-all F-rule either lowers such a count,
or lowers the total F-count. A constant step leaves all those
coordinates fixed and lowers the weight sum. Hence every step
strictly decreases this fixed finite lexicographic natural tuple.
This proves SN on every finite initial term. All nonvariable
normal forms in the pure factor are7: every other constant has
a rule and every F-root has the catch-all rule. Rewriting never
turns a nonvariable pure term into a variable. Thus any peak
from it normalizes on both sides to7; variables have no steps.
This proves CR and completeness of the first factor.

For the majority factor every step selects a proper argument,
so node count proves SN. At a common root two different majority
patterns match only when all three arguments agree, and their
contracta are identical. Disjoint occurrences commute. For a
nested step in the lone, erased argument, the root's two equal
winning arguments stay equal and the root contractum is unchanged.
For a step inside one winning argument, repeat it at its other
equal occurrence, then collapse the now equal updated pair.
The root-first branch performs that same step inside its retained
winning value. Both routes end at the same updated value.
This treats all one-step peaks, including all-equal arguments,
and proves WCR. SN and the abstract Newman argument give CR.

In the sum, change the first q-copy's1 and2 to4, the second
copy's1 and3 to5, and the third copy's2 and3 to6. The three
majority rules then yield4,5,6, respectively; keep the fourth
copy unchanged. Apply the F duplicating rule to return to four
copies of the original q. This gives the exact source cycle.
The full modularity theorem nevertheless gives CR of the sum,
since both factors were proved CR. Thus its nontermination is
not a counterexample to that theorem.
\end{proof}
\end{annexitem}

\begin{annexitem}{annex:p12-klop1987-ds-applications}{thm:p12-klop1987-ds-applications}{The actual binary and ternary factor hypotheses}
\begin{proof}
Each binary equality-test step removes a D-node and replaces
the whole redex by E, so node count decreases. Its only
nonvariable pattern position is its root; its self-overlap
has identical E-contracta. The ordinary critical-pair argument
also handles its repeated-variable inner overlaps, proving WCR.
SN and Newman give CR.

Every selector rule selects a proper argument, so node count
again proves SN. Nonvariable proper positions contain only
the irreducible constants true or false, not an if-pattern.
Root self-pairs are trivial; true and false rules cannot
overlap; and either can overlap the repeated-argument rule
only when the two branches agree, giving identical contracta.
Thus all critical pairs join, proving WCR at the full ordinary
nonlinear scope and then CR by SN. CL has CR by its separately
proved regular-confluence argument. Its full binary application
signature is disjoint from each stated factor's signature,
so the modularity theorem proves the two direct-sum applications.
Neither nonlinear factor satisfies left-linearity, and none
was assigned a conditional rewriting semantics.
\end{proof}
\end{annexitem}
